\chapter[Vorwort]{Vorwort}
\quellseite{1}

„Im Prediger Salomo heißt es einmal: ‚Viel Büchermachens ist kein Ende‘, Cap. 12, 12. War dies schon der Fall zur Zeit des alten Testaments, so ist es in noch viel größerem Maße Thatsache in unserer Zeit. Die Welt ist mit Büchern geradezu überschwemmt. Und das ist nicht durchweg ein Vortheil und ein Segen, sondern vielmehr in gar mancher Hinsicht ein rechter Schaden. Wir sehen jetzt ganz ab von den zahllosen gottlosen und verderblichen Schriften, die ein Fluch für die Menschheit sind; auch von den unzähligen andern Büchern, deren Lesen den Menschen keinen wirklichen Nutzen bringt. Es gibt auch zu viele sogenannte gute Bücher, weil diese, wie die Erfahrung zeigt, die Menschen abhalten, die besten Bücher fleißig und immer wieder zu lesen, vor allem das einzigartige Buch aller Bücher, die heilige Schrift. Darum sagte schon vor 465 Jahren ein Mann, dessen Bücher zu den besten menschlichen Schriften gehören, die je geschrieben worden sind: ‚Gern hätte ich’s gesehen, daß meine Bücher allesamt wären dahinten blieben und untergangen. Und ist unter andern Ursachen eine, daß mir grauet vor dem Exempel; denn ich wohl sehe, was Neues in der Kirche geschafft ist, da man hat außer und neben der heiligen Schrift angefangen, viel Bücher und große Bibliotheken zu sammeln, sonderlich ohn allen Unterschied allerlei Väter, Concilia und Lehrer aufzuwerfen. Damit nicht allein die edle Zeit und Studiren in der Schrift versäumt, sondern auch die reine Erkenntnuß göttliches Worts endlich verloren ist, bis die Biblia … unter der Bank im Staube vergessen ist. … Auch ist das unsere Meinung gewest, da wir die Biblia selbst zu verdeutschen anfingen, daß wir hofften, es sollt des Schreibens weniger und des Studirens und Lesens in der Schrift mehr werden. Denn auch alles andere Schreiben in und zu der Schrift, wie Johannes zu Christo, weisen soll; wie er spricht: ‚Ich muß abnehmen, dieser muß zunehmen‘ (Joh. 3, 30.), damit ein jeglicher selbst möchte aus der frischen Quelle trinken. … Nun ich’s aber ja nicht kann wehren, und man ohne meinen Dank meine Bücher will durch den Druck (mir zu kleinem Ehren) jetzt sammeln, muß ich sie die Kost und Arbeit lassen dran wagen. … Wohlan, so laß gehen in Gottes Namen, ohne daß ich freundlich bitte, wer meine Bücher zu dieser Zeit ja haben will, der lasse sie ihm bei Leibe nicht sein ein Hinderniß, die Schrift selbst zu studiren.“ So schrieb Luther im Jahre 1539 in der Vorrede über den ersten Theil seiner gesammelten deutschen Schriften\footnote{St. L. Ausg. XIV, 432 f.}.

Aber die ganzen heutigen Verhältnisse in der Welt und auch in der Kirche, die jetzt nicht näher darzulegen sind, bringen es mit sich, daß Bücher auf Bücher geschrieben und durch den Druck verbreitet werden. Wir Christen sollen
\quellseite{2}
aber bei dieser Lage der Dinge recht machen, daß wir nicht auf verkehrte Wege gezogen werden, sollen wohl prüfen, daß wir nicht mit dem Lesen unnützer Bücher die edle Zeit vergeuden, sollen namentlich nur solche religiöse Bücher lesen, die wirklich in die Schrift hineinführen, und sollen uns auch durch das beste Erbauungsbuch nicht abhalten lassen vom Studium des lieben Gotteswortes, oder die Stunden verkürzen lassen, die dem Bibellesen gewidmet sein sollten.

Was soeben von den Büchern gesagt worden ist, gilt auch von den mancherlei Zeitungen und Zeitschriften. Ja, die letzteren verdrängen fast je länger je mehr die Bücher. Es ist kaum glaublich, welch Unmaße Zeitungen in unserer Zeit gedruckt und gelesen wird. Vor uns liegt eine Zusammenstellung, aus der hervorgeht, daß die Zahl der Exemplare von Zeitungen, die nur in den bedeutenderen Ländern der Erde gedruckt werden, jährlich etwa 12,000 Millionen beträgt. In den letzten vierzig Jahren hat sich die Zahl der Zeitungen in Europa wenigstens verdoppelt und hier in America mehr als verdreifacht. Nach Zehntausenden und Hunderttausenden zählen die Häuser, in denen nicht bloß Eine, sondern zwei und noch mehr Zeitungen täglich eintreffen, des Morgens wie des Abends, Wochentags wie Sonntags. Daß diese zahllosen Zeitschriften kein reiner Segen, sondern vielmehr in gar mancher Beziehung ein Schaden sind, liegt auf der Hand. Viel zu viel Zeit wird in der Regel mit deren Lesen zugebracht, Zeit, die anders und auch mit Lesen viel nützlicher verwandt werden könnte, ganz abgesehen davon, daß die große Mehrzahl der Zeitschriften im Dienste der gottfeindlichen Welt und ihres Fürsten, des Teufels, steht, der Sünde dient und dienen will und dadurch vielen Seelen gefährlich und verderblich wird. Und je mehr nun das Zeitungswesen eine Macht und durch die ganze Weltlage freilich eine Nothwendigkeit geworden ist, desto ernstlicher ist es Christen Pflicht, dasselbe recht und nicht falsch zu gebrauchen, und auch in diesem Stücke nicht zu vergessen, daß das Wesen dieser Welt vergehet, 1 Cor. 7, 31. Nimmermehr sollen sich die Christen durch die Zeitungen die Zeit und den Geschmack für das Geistliche und Göttliche, für das theure Gotteswort in der Bibel, rauben lassen.

Auch die kirchliche Zeitschrift kann verderblich werden, und eine Menge kirchlicher Blätter fällt unter dieses Urtheil. Wir schweigen jetzt der Masse falschgläubiger Zeitschriften, die fort und fort das mörderische Seelengift der falschen Lehre austheilen. Auch ein dem Namen nach rechtgläubiges Kirchenblatt bringt keinen rechten Nutzen, wenn es lediglich ein Unterhaltungsblatt oder ein kirchliches Neuigkeitsblatt sein will und nicht die Zeiterscheinungen und die kirchlichen Ereignisse in das Licht des lauteren Gotteswortes rückt und den untrüglichen Maßstab desselben anlegt, oder wenn es, statt die Leser immer mehr in die Schrift hineinzuführen, sie von dem Worte Gottes abzieht und nicht zum Lesen desselben kommen läßt, ihnen die Schrift gewissermaßen ersetzen will. Der „Lutheraner“ hat ein anderes Ziel. Er will nicht seine Spalten füllen mit mancherlei und allerlei Nachrichten aus den verschiedenen kirchlichen Kreisen, sondern seine erste Absicht ist, in seinem Bereiche zu lehren, Gottes Wort zu lehren und die dem heiligen Schriftworte entgegenstehende falsche Lehre und alles mit der Bibel Streitende gottlose Leben zu strafen. Und was er aus dem Reiche der Kirche und der Welt

mittheilt, will er immer zugleich nach dem „Gesetz und Zeugniß“ beurtheilen und überhaupt allewege das Wort Gottes anwenden zur Lehre, zur Strafe, zur Besserung, zur Züchtigung in der Gerechtigkeit und zum Troste. Er will seine Leser im rechten biblischen Sinn des Wortes „erbauen“. Auf diese Weise hofft er, keinen seiner Leser vom fleißigen Studium des ewigen Gotteswortes abzuziehen oder die Zeit dazu zu rauben oder ihm den Geschmack daran zu nehmen, sondern vielmehr jeden Leser immer mehr in die Schrift hineinzutreiben, zum fleißigen Schriftlesen ihm rechte, heilige Lust und Liebe zu erwecken und ihn im Verständnis des Wortes der Wahrheit zu fördern und zu gründen, so daß unsere lutherischen Christen auch durch seinen Dienst immer mehr rechte Bibelchristen werden. Gott, der Herr und Hirte seiner Kirche, gebe auch im neuen Jahrgang zu dem Wollen das Vollbringen nach seinem Wohlgefallen. L. F.
